% ─── MÓDULO 7: Administración de la base de datos ────────
\newpage
\section{Administración de la Base de Datos}

Las sentencias de administración de bases de datos pertenecen al sublenguaje
\textbf{DDL}. Operan sobre el contenedor de más alto nivel en MySQL: la
\textbf{base de datos} (también llamada \textit{schema}), que agrupa tablas,
vistas, procedimientos y demás objetos.

% ─────────────────────────────────────────────────────────
\subsection{\texttt{CREATE DATABASE}}

\begin{definicion}{CREATE DATABASE}
  Crea una nueva base de datos vacía en el servidor. Permite especificar el
  \textbf{juego de caracteres} (\textit{character set}) y la
  \textbf{colación} (\textit{collation}) que determinarán cómo se almacenan
  y comparan las cadenas de texto en todas sus tablas.
\end{definicion}

\begin{lstlisting}[style=sql, caption={Sintaxis de CREATE DATABASE}]
CREATE DATABASE [IF NOT EXISTS] nombre_bd
    [CHARACTER SET  charset_name]
    [COLLATE        collation_name];
\end{lstlisting}

\begin{ejemplo}
  Crear una base de datos para una aplicación en español con soporte
  completo de Unicode:

\begin{lstlisting}[style=sql, caption={Base de datos con utf8mb4 y collation en español}]
CREATE DATABASE IF NOT EXISTS tienda_online
    CHARACTER SET  utf8mb4
    COLLATE        utf8mb4_unicode_ci;

-- Seleccionar la base de datos activa
USE tienda_online;
\end{lstlisting}
\end{ejemplo}

\textbf{Juegos de caracteres y colaciones más comunes:}

\begin{table}[H]
  \centering
  \caption{Character sets y collations habituales en MySQL}
  \renewcommand{\arraystretch}{1.3}
  \rowcolors{2}{gray!10}{white}
  \begin{tabularx}{\textwidth}{>{\ttfamily}l >{\ttfamily}l X}
    \toprule
    \textbf{charset}  & \textbf{collation}          & \textbf{Uso recomendado}                        \\
    \midrule
    utf8mb4  & utf8mb4\_unicode\_ci  & Unicode completo; soporta emojis; recomendado   \\
    utf8mb4  & utf8mb4\_0900\_ai\_ci & Unicode 9.0; sin acento/mayúsculas; MySQL 8.0+  \\
    utf8mb4  & utf8mb4\_bin          & Comparación binaria; distingue mayúsculas       \\
    latin1   & latin1\_swedish\_ci   & Legado; solo caracteres ISO-8859-1              \\
    \bottomrule
  \end{tabularx}
\end{table}

\begin{nota}
  Usar \texttt{utf8} (sin \texttt{mb4}) en MySQL es un error frecuente: ese
  alias solo soporta hasta 3 bytes por carácter y excluye emojis y algunos
  ideogramas CJK. La elección correcta para proyectos nuevos es siempre
  \texttt{utf8mb4}.
\end{nota}

\begin{lstlisting}[style=sql, caption={Comandos de exploración de bases de datos}]
SHOW DATABASES;                    -- listar todas las bases de datos
SHOW CREATE DATABASE tienda_online; -- ver la definición completa
SELECT DATABASE();                  -- ver la base de datos activa
\end{lstlisting}

% ─────────────────────────────────────────────────────────
\subsection{\texttt{DROP DATABASE}}

\begin{definicion}{DROP DATABASE}
  Elimina la base de datos y \textbf{todos los objetos que contiene}
  (tablas, vistas, procedimientos, datos) de forma permanente e
  irreversible. No existe confirmación ni papelera de reciclaje.
\end{definicion}

\begin{lstlisting}[style=sql, caption={Sintaxis de DROP DATABASE}]
DROP DATABASE [IF EXISTS] nombre_bd;
\end{lstlisting}

\begin{ejemplo}
\begin{lstlisting}[style=sql, caption={DROP DATABASE con y sin guarda}]
-- Genera error si la base de datos no existe
DROP DATABASE tienda_online;

-- No genera error si no existe (útil en scripts de despliegue)
DROP DATABASE IF EXISTS tienda_online;
\end{lstlisting}
\end{ejemplo}

\begin{nota}
  \textbf{Precauciones antes de ejecutar DROP DATABASE:}
  \begin{enumerate}
    \item Verificar que se tiene un respaldo (\textit{backup}) actualizado.
    \item Confirmar que ninguna aplicación activa tiene conexiones abiertas
          a esa base de datos.
    \item Comprobar el nombre con \texttt{SHOW DATABASES} antes de ejecutar.
    \item En producción, revocar el privilegio \texttt{DROP} a usuarios de
          aplicación para evitar ejecuciones accidentales.
  \end{enumerate}
\end{nota}

% ─── MÓDULO 13: Conceptos avanzados ──────────────────────
\newpage
\section{Conceptos Avanzados}

Los objetos descritos en este módulo permiten construir bases de datos más
eficientes, seguras y mantenibles. Aunque no son imprescindibles para
consultas básicas, son herramientas fundamentales en entornos de producción.

% ─────────────────────────────────────────────────────────
\subsection{\texttt{INDEX}}

\begin{definicion}{Índice}
  Estructura de datos auxiliar que el motor mantiene sincronizada con la
  tabla para acelerar las operaciones de búsqueda, ordenamiento y unión.
  Mejora el rendimiento de \texttt{SELECT} a cambio de un coste adicional
  en \texttt{INSERT}, \texttt{UPDATE} y \texttt{DELETE}, y de espacio en
  disco.
\end{definicion}

\textbf{Tipos de índice en MySQL:}

\begin{table}[H]
  \centering
  \caption{Tipos de índice disponibles en MySQL/InnoDB}
  \renewcommand{\arraystretch}{1.3}
  \rowcolors{2}{gray!10}{white}
  \begin{tabularx}{\textwidth}{>{\ttfamily}l l X}
    \toprule
    \textbf{Tipo}   & \textbf{Motor}       & \textbf{Uso recomendado}                              \\
    \midrule
    B-TREE          & InnoDB, MyISAM       & Búsquedas exactas, rangos, \texttt{ORDER BY}          \\
    HASH            & Memory               & Solo igualdad exacta; no soporta rangos               \\
    FULLTEXT        & InnoDB (5.6+)        & Búsqueda de texto libre en columnas \texttt{TEXT}     \\
    SPATIAL         & InnoDB, MyISAM       & Datos geoespaciales (\texttt{GEOMETRY})               \\
    \bottomrule
  \end{tabularx}
\end{table}

\begin{lstlisting}[style=sql, caption={Creación y eliminación de índices}]
-- Índice simple en una columna
CREATE INDEX idx_apellido ON clientes(apellido);

-- Índice compuesto (el orden de columnas importa)
CREATE INDEX idx_cat_precio ON productos(categoria_id, precio);

-- Índice UNIQUE (equivale a restricción UNIQUE)
CREATE UNIQUE INDEX uq_correo ON usuarios(correo);

-- Índice FULLTEXT para búsqueda de texto
CREATE FULLTEXT INDEX ft_descripcion ON productos(descripcion);

-- Eliminar un índice
DROP INDEX idx_apellido ON clientes;
\end{lstlisting}

\begin{lstlisting}[style=sql, caption={Uso del índice FULLTEXT con MATCH ... AGAINST}]
SELECT nombre, descripcion
FROM   productos
WHERE  MATCH(descripcion) AGAINST('teclado mecánico' IN BOOLEAN MODE);
\end{lstlisting}

\begin{nota}
  \textbf{Cuándo crear un índice:}
  \begin{itemize}
    \item Columnas que aparecen frecuentemente en \texttt{WHERE}, \texttt{JOIN ON},
          \texttt{ORDER BY} o \texttt{GROUP BY}.
    \item Columnas con alta cardinalidad (muchos valores distintos).
  \end{itemize}
  \textbf{Cuándo evitarlo:}
  \begin{itemize}
    \item Tablas pequeñas (el escaneo completo es más rápido).
    \item Columnas con muy baja cardinalidad (\texttt{activo}, \texttt{genero}).
    \item Tablas con escrituras masivas y pocas lecturas.
  \end{itemize}
\end{nota}

% ─────────────────────────────────────────────────────────
\subsection{\texttt{TRIGGER}}

\begin{definicion}{TRIGGER}
  Bloque de código SQL que se ejecuta \textbf{automáticamente} en respuesta
  a un evento DML (\texttt{INSERT}, \texttt{UPDATE} o \texttt{DELETE}) sobre
  una tabla, antes (\texttt{BEFORE}) o después (\texttt{AFTER}) de que el
  cambio se aplique.
\end{definicion}

\begin{lstlisting}[style=sql, caption={Sintaxis general de un TRIGGER}]
CREATE TRIGGER nombre_trigger
    { BEFORE | AFTER } { INSERT | UPDATE | DELETE }
    ON nombre_tabla
    FOR EACH ROW
BEGIN
    -- cuerpo del trigger
    -- NEW.columna  -> valor nuevo (INSERT / UPDATE)
    -- OLD.columna  -> valor anterior (UPDATE / DELETE)
END;
\end{lstlisting}

\begin{ejemplo}
  Registrar en una tabla de auditoría cada cambio de precio en
  \texttt{productos}:

\begin{lstlisting}[style=sql, caption={TRIGGER de auditoría de cambios de precio}]
CREATE TABLE auditoria_precios (
    id          INT UNSIGNED AUTO_INCREMENT PRIMARY KEY,
    producto_id INT UNSIGNED NOT NULL,
    precio_ant  DECIMAL(10,2),
    precio_new  DECIMAL(10,2),
    cambiado_en DATETIME NOT NULL DEFAULT CURRENT_TIMESTAMP
);

DELIMITER $$
CREATE TRIGGER trg_auditoria_precio
    AFTER UPDATE ON productos
    FOR EACH ROW
BEGIN
    IF OLD.precio <> NEW.precio THEN
        INSERT INTO auditoria_precios
               (producto_id, precio_ant, precio_new)
        VALUES (NEW.producto_id, OLD.precio, NEW.precio);
    END IF;
END$$
DELIMITER ;
\end{lstlisting}
\end{ejemplo}

\begin{nota}
  Los triggers dificultan la depuración y pueden crear efectos secundarios
  inesperados si no se documentan. Úsarlos con moderación; para lógica
  compleja, preferir procedimientos almacenados invocados explícitamente.
\end{nota}

% ─────────────────────────────────────────────────────────
\subsection{\texttt{VIEW}}

\begin{definicion}{Vista (\textit{View})}
  Consulta \texttt{SELECT} almacenada con nombre propio que puede usarse
  como si fuera una tabla. No duplica datos: cada vez que se consulta, el
  motor ejecuta la consulta subyacente. Simplifica consultas complejas,
  encapsula lógica y puede restringir el acceso a columnas sensibles.
\end{definicion}

\begin{lstlisting}[style=sql, caption={Creación y consulta de una vista}]
-- Crear la vista
CREATE VIEW vista_pedidos_activos AS
    SELECT p.pedido_id,
           c.nombre          AS cliente,
           p.total,
           p.fecha
    FROM   pedidos    AS p
    INNER JOIN clientes AS c ON p.cliente_id = c.cliente_id
    WHERE  p.estado = 'activo';

-- Usar la vista exactamente como una tabla
SELECT cliente, SUM(total) AS total_compras
FROM   vista_pedidos_activos
GROUP  BY cliente
ORDER  BY total_compras DESC;

-- Reemplazar o eliminar la vista
CREATE OR REPLACE VIEW vista_pedidos_activos AS ...;
DROP VIEW IF EXISTS vista_pedidos_activos;
\end{lstlisting}

\begin{nota}
  Las vistas en MySQL no almacenan datos en caché por defecto. Para
  escenarios de alto rendimiento con consultas agregadas costosas,
  considerar \textbf{vistas materializadas} (disponibles en PostgreSQL) o
  tablas de resumen actualizadas por eventos programados.
\end{nota}

% ─────────────────────────────────────────────────────────
\subsection{\texttt{STORED PROCEDURE}}

\begin{definicion}{Procedimiento almacenado}
  Bloque de código SQL con nombre, almacenado en el servidor y ejecutable
  mediante \texttt{CALL}. Acepta parámetros de entrada (\texttt{IN}),
  salida (\texttt{OUT}) o ambos (\texttt{INOUT}), y permite estructuras de
  control (\texttt{IF}, \texttt{WHILE}, \texttt{LOOP}, cursores).
\end{definicion}

\begin{lstlisting}[style=sql, caption={Procedimiento con parámetros IN y OUT}]
DELIMITER $$
CREATE PROCEDURE sp_resumen_cliente(
    IN  p_cliente_id INT UNSIGNED,
    OUT p_total_pedidos INT,
    OUT p_monto_total   DECIMAL(12,2)
)
BEGIN
    SELECT COUNT(*),
           COALESCE(SUM(total), 0)
    INTO   p_total_pedidos,
           p_monto_total
    FROM   pedidos
    WHERE  cliente_id = p_cliente_id;
END$$
DELIMITER ;

-- Invocar el procedimiento
CALL sp_resumen_cliente(7, @num_pedidos, @monto);
SELECT @num_pedidos AS pedidos, @monto AS monto_total;
\end{lstlisting}

\begin{ejemplo}
  Procedimiento con flujo de control para aplicar descuentos escalonados:

\begin{lstlisting}[style=sql, caption={Stored procedure con IF / ELSEIF}]
DELIMITER $$
CREATE PROCEDURE sp_aplicar_descuento(
    IN p_producto_id INT UNSIGNED,
    IN p_porcentaje  DECIMAL(5,2)
)
BEGIN
    IF p_porcentaje <= 0 OR p_porcentaje > 80 THEN
        SIGNAL SQLSTATE '45000'
            SET MESSAGE_TEXT = 'Porcentaje fuera de rango permitido';
    ELSEIF p_porcentaje <= 20 THEN
        UPDATE productos
        SET    precio = precio * (1 - p_porcentaje / 100)
        WHERE  producto_id = p_producto_id;
    ELSE
        -- Descuento mayor requiere aprobación: solo registrar solicitud
        INSERT INTO solicitudes_descuento(producto_id, porcentaje)
        VALUES (p_producto_id, p_porcentaje);
    END IF;
END$$
DELIMITER ;
\end{lstlisting}
\end{ejemplo}

% ─────────────────────────────────────────────────────────
\subsection{Transacciones}

\begin{definicion}{Transacción}
  Unidad lógica de trabajo que agrupa una o más operaciones DML. Garantiza
  las propiedades \textbf{ACID}: Atomicidad, Consistencia, Aislamiento y
  Durabilidad. Si alguna operación falla, todo el grupo puede deshacerse.
\end{definicion}

\begin{table}[H]
  \centering
  \caption{Comandos de control de transacciones}
  \renewcommand{\arraystretch}{1.3}
  \rowcolors{2}{gray!10}{white}
  \begin{tabularx}{\textwidth}{>{\ttfamily}l X}
    \toprule
    \textbf{Comando}          & \textbf{Efecto}                                               \\
    \midrule
    START TRANSACTION / BEGIN & Inicia una transacción explícita                              \\
    COMMIT                    & Confirma y persiste todos los cambios                         \\
    ROLLBACK                  & Deshace todos los cambios desde el último COMMIT              \\
    SAVEPOINT nombre          & Crea un punto de restauración parcial dentro de la transacción\\
    ROLLBACK TO nombre        & Deshace hasta el savepoint indicado (sin terminar la transacción)\\
    RELEASE SAVEPOINT nombre  & Elimina un savepoint sin hacer rollback                       \\
    \bottomrule
  \end{tabularx}
\end{table}

\begin{lstlisting}[style=sql, caption={Transacción con SAVEPOINT y manejo de errores}]
START TRANSACTION;

    INSERT INTO pedidos (cliente_id, total, fecha)
    VALUES (12, 1500.00, CURDATE());

    SAVEPOINT sp_pedido_insertado;

    INSERT INTO detalle_pedido (pedido_id, producto_id, cantidad)
    VALUES (LAST_INSERT_ID(), 5, 3);

    -- Si este UPDATE falla, se puede revertir solo el detalle
    UPDATE productos SET stock = stock - 3 WHERE producto_id = 5;

    -- Todo correcto: confirmar
    COMMIT;

-- En caso de error detectado en la aplicación:
-- ROLLBACK TO sp_pedido_insertado;  -- deshace solo desde el savepoint
-- ROLLBACK;                          -- deshace toda la transacción
\end{lstlisting}

\begin{nota}
  MySQL opera en modo \texttt{AUTOCOMMIT = 1} por defecto: cada sentencia
  DML es una transacción implícita que se confirma de inmediato.
  \texttt{START TRANSACTION} suspende ese modo hasta que se ejecute
  \texttt{COMMIT} o \texttt{ROLLBACK}.
\end{nota}

% ─────────────────────────────────────────────────────────
\subsection{Concurrencia}

\begin{definicion}{Concurrencia}
  Capacidad del motor para atender \textbf{múltiples transacciones
  simultáneas} sin que interfieran entre sí. Se gestiona mediante
  \textbf{bloqueos} (\textit{locks}) y \textbf{niveles de aislamiento}
  (\textit{isolation levels}).
\end{definicion}

\textbf{Niveles de aislamiento en SQL estándar:}

\begin{table}[H]
  \centering
  \caption{Niveles de aislamiento y fenómenos que previenen}
  \renewcommand{\arraystretch}{1.3}
  \rowcolors{2}{gray!10}{white}
  \begin{tabularx}{\textwidth}{l X c c c}
    \toprule
    \textbf{Nivel}      & \textbf{Descripción resumida}
      & \textbf{\textit{Dirty read}}
      & \textbf{\textit{Non-rep. read}}
      & \textbf{\textit{Phantom}} \\
    \midrule
    READ UNCOMMITTED & Ve cambios no confirmados de otras transacciones   & Sí & Sí & Sí \\
    READ COMMITTED   & Solo ve cambios ya confirmados                     & No & Sí & Sí \\
    REPEATABLE READ  & Garantiza lecturas consistentes (defecto InnoDB)   & No & No & Sí \\
    SERIALIZABLE     & Ejecución secuencial total; máximo aislamiento      & No & No & No \\
    \bottomrule
  \end{tabularx}
\end{table}

\begin{lstlisting}[style=sql, caption={Configurar y consultar el nivel de aislamiento}]
-- Ver el nivel actual
SELECT @@transaction_isolation;

-- Cambiar para la sesión actual
SET SESSION TRANSACTION ISOLATION LEVEL REPEATABLE READ;

-- Cambiar globalmente (requiere privilegio SUPER)
SET GLOBAL  TRANSACTION ISOLATION LEVEL READ COMMITTED;
\end{lstlisting}

\begin{nota}
  \textbf{Deadlock:} dos transacciones se bloquean mutuamente esperando
  recursos que la otra retiene. InnoDB detecta el ciclo automáticamente y
  aborta la transacción de menor coste. La aplicación debe estar preparada
  para reintentar la operación al recibir el error \texttt{ER\_LOCK\_DEADLOCK}
  (\texttt{1213}). Para minimizarlos: acceder siempre a las tablas en el
  mismo orden y mantener las transacciones lo más cortas posible.
\end{nota}

% ─── MÓDULO 14: Conexión desde código ────────────────────
\newpage
\section{Conexión desde Código}

Integrar una base de datos en una aplicación requiere un \textbf{conector}
o \textit{driver}: biblioteca que traduce las llamadas del lenguaje de
programación al protocolo de red del motor SQL.

% ─────────────────────────────────────────────────────────
\subsection{Conectores y Drivers}

\begin{table}[H]
  \centering
  \caption{Conectores oficiales y populares por lenguaje}
  \renewcommand{\arraystretch}{1.3}
  \rowcolors{2}{gray!10}{white}
  \begin{tabularx}{\textwidth}{l >{\ttfamily}l X}
    \toprule
    \textbf{Lenguaje} & \textbf{Librería / Driver}  & \textbf{Instalación / Referencia}          \\
    \midrule
    Python   & mysql-connector-python & \texttt{pip install mysql-connector-python} \\
    Python   & PyMySQL                & \texttt{pip install pymysql}                \\
    Python   & SQLAlchemy (ORM)       & \texttt{pip install sqlalchemy}             \\
    Node.js  & mysql2                 & \texttt{npm install mysql2}                 \\
    Node.js  & Prisma (ORM)           & \texttt{npm install prisma}                 \\
    Java     & MySQL Connector/J      & Dependencia Maven: \texttt{mysql:mysql-connector-java} \\
    PHP      & PDO\_MySQL             & Incluido en PHP; activar extensión \texttt{pdo\_mysql} \\
    Go       & go-sql-driver/mysql    & \texttt{go get github.com/go-sql-driver/mysql} \\
    \bottomrule
  \end{tabularx}
\end{table}

\begin{ejemplo}
  Conexión y consulta básica con Python (\texttt{mysql-connector-python}):

\begin{lstlisting}[language=Python, basicstyle=\ttfamily\small,
    keywordstyle=\bfseries\color{NavyBlue},
    commentstyle=\itshape\color{Gray},
    backgroundcolor=\color{gray!6}, frame=single,
    caption={Conexión a MySQL desde Python}]
import mysql.connector

conn = mysql.connector.connect(
    host     = "localhost",
    user     = "app_user",
    password = "s3cr3to",
    database = "tienda_online"
)

cursor = conn.cursor(dictionary=True)
cursor.execute("SELECT producto_id, nombre, precio FROM productos LIMIT 5")

for fila in cursor.fetchall():
    print(fila["nombre"], fila["precio"])

cursor.close()
conn.close()
\end{lstlisting}
\end{ejemplo}

\begin{ejemplo}
  Conexión con Node.js usando \texttt{mysql2} y \texttt{async/await}:

\begin{lstlisting}[language=Java, basicstyle=\ttfamily\small,
    keywordstyle=\bfseries\color{NavyBlue},
    commentstyle=\itshape\color{Gray},
    backgroundcolor=\color{gray!6}, frame=single,
    caption={Conexión a MySQL desde Node.js}]
const mysql = require('mysql2/promise');

const conn = await mysql.createConnection({
    host:     'localhost',
    user:     'app_user',
    password: 's3cr3to',
    database: 'tienda_online'
});

const [rows] = await conn.execute(
    'SELECT nombre, precio FROM productos WHERE categoria_id = ?',
    [3]   // parámetro enlazado (previene SQL Injection)
);

console.table(rows);
await conn.end();
\end{lstlisting}
\end{ejemplo}

% ─────────────────────────────────────────────────────────
\subsection{SQL Injection}

\begin{definicion}{SQL Injection}
  Vulnerabilidad que permite a un atacante \textbf{inyectar código SQL
  arbitrario} en una consulta, manipulando su lógica para extraer,
  modificar o destruir datos, o para eludir la autenticación.
\end{definicion}

\begin{ejemplo}
  Código vulnerable: concatenación directa de entrada del usuario.

\begin{lstlisting}[language=Python, basicstyle=\ttfamily\small,
    keywordstyle=\bfseries\color{NavyBlue},
    commentstyle=\itshape\color{Gray},
    backgroundcolor=\color{gray!6}, frame=single,
    caption={Código VULNERABLE a SQL Injection}]
usuario = input("Usuario: ")   # atacante ingresa:  ' OR '1'='1
query = "SELECT * FROM usuarios WHERE username = '" + usuario + "'"
# Consulta resultante:
# SELECT * FROM usuarios WHERE username = '' OR '1'='1'
# -> devuelve TODOS los usuarios sin contraseña
cursor.execute(query)
\end{lstlisting}
\end{ejemplo}

\begin{ejemplo}
  Código seguro: consultas parametrizadas (\textit{prepared statements}).

\begin{lstlisting}[language=Python, basicstyle=\ttfamily\small,
    keywordstyle=\bfseries\color{NavyBlue},
    commentstyle=\itshape\color{Gray},
    backgroundcolor=\color{gray!6}, frame=single,
    caption={Prevención con consultas parametrizadas}]
usuario = input("Usuario: ")
# El driver escapa automáticamente el valor
cursor.execute(
    "SELECT * FROM usuarios WHERE username = %s",
    (usuario,)   # parámetro enlazado, nunca concatenado
)
\end{lstlisting}
\end{ejemplo}

\textbf{Medidas de prevención complementarias:}

\begin{itemize}
  \item Usar \textbf{siempre} consultas parametrizadas o ORMs que las
        generen internamente.
  \item Aplicar el \textbf{principio de mínimo privilegio}: el usuario de
        la aplicación solo debe tener los permisos estrictamente necesarios
        (\texttt{SELECT}, \texttt{INSERT}, etc.), nunca \texttt{DROP} o
        \texttt{GRANT}.
  \item Validar y sanitizar la entrada en la capa de aplicación antes de
        pasarla al driver.
  \item Evitar exponer mensajes de error detallados del motor en respuestas
        de la API.
\end{itemize}

% ─── MÓDULO 15: Otros clientes y sistemas ────────────────
\newpage
\section{Otros Clientes SQL}

Existen múltiples herramientas gráficas y de línea de comandos para
interactuar con bases de datos SQL sin necesidad de escribir directamente
en la terminal.

\begin{table}[H]
  \centering
  \caption{Clientes SQL populares}
  \renewcommand{\arraystretch}{1.3}
  \rowcolors{2}{gray!10}{white}
  \begin{tabularx}{\textwidth}{l l X}
    \toprule
    \textbf{Cliente}  & \textbf{Plataforma}          & \textbf{Características destacadas}                  \\
    \midrule
    DBeaver           & Windows, macOS, Linux        & Gratuito y de código abierto; soporta más de 80 motores \\
    TablePlus         & macOS, Windows, Linux, iOS   & Interfaz nativa moderna; licencia de pago con prueba gratuita \\
    DataGrip          & Windows, macOS, Linux        & IDE de JetBrains; refactorización SQL, inspecciones inteligentes \\
    pgAdmin           & Web / escritorio             & Cliente oficial de PostgreSQL; gestión completa del servidor \\
    Adminer           & Web (un solo archivo PHP)    & Ligero, sin instalación; útil en servidores remotos  \\
    MySQL Workbench   & Windows, macOS, Linux        & Cliente oficial de MySQL; diseño ER visual incluido   \\
    \bottomrule
  \end{tabularx}
\end{table}

% ─────────────────────────────────────────────────────────
\section{PostgreSQL}

PostgreSQL es un sistema de gestión de bases de datos relacionales
\textbf{de código abierto}, con más de 35 años de desarrollo activo,
reconocido por su robustez, extensibilidad y cumplimiento estricto del
estándar SQL.

\textbf{Diferencias clave con MySQL:}

\begin{table}[H]
  \centering
  \caption{MySQL vs.\ PostgreSQL — diferencias principales}
  \renewcommand{\arraystretch}{1.3}
  \rowcolors{2}{gray!10}{white}
  \begin{tabularx}{\textwidth}{l X X}
    \toprule
    \textbf{Aspecto}          & \textbf{MySQL}                          & \textbf{PostgreSQL}                        \\
    \midrule
    Licencia                  & GPL / comercial (Oracle)                & PostgreSQL License (MIT-like)              \\
    Modelo MVCC               & Solo en InnoDB                          & Nativo en todo el motor                    \\
    Tipos de datos avanzados  & Limitados                               & JSON/JSONB, Arrays, HSTORE, UUID nativo    \\
    Vistas materializadas     & No nativo                               & Sí, con \texttt{REFRESH MATERIALIZED VIEW} \\
    Full-text search          & Básico                                  & Avanzado; soporta diccionarios y ranking   \\
    Extensiones               & Plugins limitados                       & PostGIS, pgcrypto, TimescaleDB, etc.       \\
    Procedimientos almacenados& SQL / MySQL dialect                     & PL/pgSQL, PL/Python, PL/Perl, etc.        \\
    Secuencias                & \texttt{AUTO\_INCREMENT}                & \texttt{SERIAL} / \texttt{GENERATED ALWAYS AS IDENTITY} \\
    \bottomrule
  \end{tabularx}
\end{table}

\begin{nota}
  PostgreSQL es generalmente preferido para aplicaciones que requieren
  consultas complejas, integridad de datos estricta o tipos de datos
  avanzados. MySQL sigue siendo muy popular en aplicaciones web de alto
  tráfico gracias a su velocidad en lecturas simples y su amplio soporte
  en servicios de alojamiento.
\end{nota}

% ─────────────────────────────────────────────────────────
\section{Despliegue en la Nube}

Los servicios gestionados (\textit{Database as a Service}, DBaaS) eliminan
la carga operativa de instalar, parchear y respaldar el motor de base de
datos.

\begin{table}[H]
  \centering
  \caption{Servicios de base de datos en la nube}
  \renewcommand{\arraystretch}{1.3}
  \rowcolors{2}{gray!10}{white}
  \begin{tabularx}{\textwidth}{l l X}
    \toprule
    \textbf{Servicio}     & \textbf{Motor(es)}           & \textbf{Características destacadas}                     \\
    \midrule
    AWS RDS               & MySQL, PostgreSQL, MariaDB, Oracle, SQL Server
                                                          & Backups automáticos, Multi-AZ, réplicas de lectura      \\
    AWS Aurora            & MySQL / PostgreSQL compatible& Rendimiento hasta 5$\times$ sobre RDS; serverless disponible \\
    Google Cloud SQL      & MySQL, PostgreSQL, SQL Server& Integración con BigQuery y Vertex AI                    \\
    Azure Database        & MySQL, PostgreSQL            & Alta disponibilidad con zona de disponibilidad flexible  \\
    PlanetScale           & MySQL (Vitess)               & Branching de esquemas; escala horizontal automática      \\
    Railway               & MySQL, PostgreSQL, MongoDB   & Despliegue en segundos; ideal para prototipos            \\
    Supabase              & PostgreSQL                   & BaaS open-source; autenticación y API REST incluidos     \\
    Neon                  & PostgreSQL                   & Serverless; branching por pull-request                   \\
    \bottomrule
  \end{tabularx}
\end{table}

% ─────────────────────────────────────────────────────────
\section{Próximos Pasos}

\subsection*{Caminos de especialización}

\begin{itemize}
  \item \textbf{DBA (\textit{Database Administrator}):} administración de
        servidores, alta disponibilidad, replicación, \textit{tuning} de
        rendimiento, backups y recuperación ante desastres.
  \item \textbf{Data Engineer:} pipelines de datos (ETL/ELT), almacenes de
        datos (\textit{data warehouses}), herramientas como dbt, Apache
        Spark, Airflow y servicios como BigQuery o Redshift.
  \item \textbf{Backend Developer:} integración de bases de datos en APIs
        REST y GraphQL, ORM avanzados, patrones de diseño de esquemas,
        caché con Redis.
  \item \textbf{Data Analyst / BI:} consultas analíticas avanzadas,
        funciones de ventana, herramientas de visualización (Metabase,
        Tableau, Power BI).
\end{itemize}

\subsection*{Recursos recomendados}

\begin{itemize}
  \item \textbf{Documentación oficial:}
        \href{https://dev.mysql.com/doc/}{dev.mysql.com/doc} y
        \href{https://www.postgresql.org/docs/}{postgresql.org/docs}.
  \item \textbf{Práctica interactiva:}
        \href{https://sqlzoo.net}{SQLZoo},
        \href{https://www.hackerrank.com/domains/sql}{HackerRank SQL} y
        \href{https://leetcode.com/problemset/database/}{LeetCode Database}.
  \item \textbf{Libros de referencia:}
        \textit{Learning MySQL} (Seyed M.M. \& Jay Pipes),
        \textit{PostgreSQL: Up and Running} (Regina O. Obe \& Leo S. Hsu),
        \textit{Designing Data-Intensive Applications} (Martin Kleppmann).
  \item \textbf{Modelado de datos:}
        \href{https://dbdiagram.io}{dbdiagram.io} para diseño visual de
        esquemas con notación DBML.
\end{itemize}

\printbibliography[heading=bibintoc, title={Referencias}]